%=============================================================================!
% E.2  PRECISION AND DEVICES                                                  !
%=============================================================================!
\section{Precision and devices}
\label{sec:env-precision}

A computer stores a real number in floating point: a fixed number of
significant binary digits and an exponent. Double precision
(\texttt{float64}) keeps 53 binary digits and single precision
(\texttt{float32}) keeps 24, so the gap between a number $x$ and the next
representable one is about $2^{-52}\lvert x\rvert$ in the first case and
$2^{-23}\lvert x\rvert$ in the second. The size of that gap decides which
changes in a quantity can be seen at all. The potential energy of a
simulation cell is a sum over every atom and grows with their number;
for a total energy of \EnvTotalEnergy, the gap is \EnvGapDouble{} in
double precision and \EnvGapSingle{} in single precision. Conservation of
energy, the test of an integrator in Chapter~4, is judged by changes in
the total energy far smaller than the energy itself, and a change smaller
than the gap is lost in rounding. Integrators and every test of energy
conservation in this book therefore run in double precision.

Training a neural network is usually done in single precision, which
halves the memory of every stored number and is the arithmetic for which
graphics processors (GPUs) are built; Chapter~14 checks for our own models
that the choice does not change the result. PyTorch reaches the GPU of a recent
Mac through its Metal backend, \texttt{mps}, which supports single but not
double precision: asking it for a \texttt{float64} tensor raises
\EnvMpsError. The division of labour in this book follows from that.
Integration and conservation studies run on the CPU in double precision;
network training may run on the \texttt{mps} device or a CUDA GPU in
single precision, and where the book uses a GPU it measures the speed-up
against the CPU rather than assuming one.
